% PAQUETES DEL PROYECTO

% Configuración de codificación y lenguaje
\usepackage[utf8]{inputenc} % Caracteres UTF-8
\usepackage[T1]{fontenc}
\usepackage[spanish]{babel} % Idioma español

% Letra
\usepackage{libertinus}
\usepackage{libertinust1math}

\usepackage{FiraMono}
\renewcommand{\ttdefault}{fvm}
\DeclareTextFontCommand{\texttt}{\footnotesize\ttfamily}

% Referencias bibliográficas
\usepackage[backend=biber, style=ieee, language=spanish]{biblatex}
\addbibresource{references.bib}
\usepackage{csquotes} % Comillas españolas en las citas

% Formato de documentos y márgenes
\usepackage[left=18mm,right=18mm,top=21mm,bottom=21mm]{geometry}
\usepackage{fancyhdr} % Encabezados y pies de página
\flushbottom % Todas las páginas tengan la misma altura
\usepackage{lastpage} % Última página

%Manejo de fechas
\usepackage[es]{datetime2}

% Leyendas de imágenes/tablas
\usepackage{caption}
\DeclareCaptionLabelSeparator{midsep}{\hspace{0.5em}}
\captionsetup[figure]{
  labelfont=bf,
  labelsep=midsep
}
\captionsetup[table]{
  labelfont=bf,
  labelsep=midsep
}
\usepackage{subcaption}
\captionsetup[table]{name=Tabla}

% Figuras
\usepackage{pgfplots}
\pgfplotsset{compat=1.18}

% Ubicación de elementos flotantes como figuras y tablas
\usepackage{float}
\renewcommand{\topfraction}{0.9} % Máx. parte superior ocupada por floats
\renewcommand{\bottomfraction}{0.8} % Máx. parte inferior
\renewcommand{\textfraction}{0.1} % Mínimo texto que debe quedar
\renewcommand{\floatpagefraction}{0.8} % Requerido para página solo con floats

% Paquetes de matemática
\usepackage{amsmath, amsthm, amssymb}
\usepackage{bm} % bold math
\usepackage{mathtools}
\usepackage{nicefrac}
\usepackage{commath}
\decimalpoint % Usa punto decimal en lugar de coma.

% Circuitos
\usepackage[american, siunitx]{circuitikz}

% Unidades del SI
\usepackage{siunitx}

% Paquete para insertar imágenes y gráficos
\usepackage{import, graphicx}

% Tablas
\usepackage{tabularray}
\usepackage{booktabs}
\usepackage{multirow}
\usepackage{multicol}
\newcolumntype{M}[1]{>{\centering\arraybackslash}m{#1}}

% Tikz
\usepackage{graphics,tikz}
\usetikzlibrary{babel}

% Listas numeradas y viñetas
\usepackage{enumitem}
\setlist[itemize]{label=\textbullet}

% Letra de títulos
\usepackage{titlesec}

% Section con tamaño de 15pt
\titleformat{\section}
  {\normalfont\fontsize{15}{16}\selectfont\bfseries}
  {\thesection}
  {1em}
  {}

% Subsection con tamaño de 12pt
\titleformat{\subsection}
  {\normalfont\fontsize{12}{14}\selectfont\bfseries}
  {\thesubsection}
  {1em}
  {}

% Subsubsection con tamaño de 10pt
\titleformat{\subsubsection}
  {\normalfont\fontsize{10}{12}\selectfont\bfseries}
  {\thesubsubsection}
  {1em}
  {}

% Paragraph
\titleformat{\paragraph}
  {\normalfont\fontsize{10}{12}\selectfont\bfseries}
  {\theparagraph}
  {0.5em}
  {}

% Entorno inciso
\usepackage[most]{tcolorbox}
% Contador
\newcounter{inciso}[subsection]
\renewcommand{\theinciso}{\alph{inciso}}
% Estilo del inciso
\newtcolorbox{incisobox}{
  enhanced,
  sharp corners,
  colback=white,
  colframe=black,
  boxrule=0pt,
  frame hidden,
  borderline west={2pt}{0pt}{black},
  left=10pt,
  right=0pt,
  top=4pt,
  bottom=4pt,
}
% Comando
\newcommand{\inciso}[1]{%
  \refstepcounter{inciso}%
  \addcontentsline{toc}{subsubsection}{Inciso \theinciso}
  \begin{incisobox}
    \textbf{\theinciso)} #1
  \end{incisobox}
}

% Manejo de enlaces
\usepackage[colorlinks=true, linkcolor=black, citecolor=blue, urlcolor=blue]{hyperref}

% Compatibilidad con entorno IEEEkeywords (venía del formato IEEEtran)
\newenvironment{IEEEkeywords}{\par\smallskip\noindent\textbf{Palabras clave:} \small}{\par}
\renewcommand{\UrlFont}{\footnotesize\ttfamily}
% - Activa los enlaces en el documento
% - linkcolor=black (enlaces internos en negro)
% - citecolor=blue (citas en azul)
% - urlcolor=blue (url en azul)
